\section{Supporting Experiments (from Main Text)}
\label{sec:supporting-experiments}

The following results extend EVQ beyond the primary MHA setting. Each is scoped to a specific question and labeled with its seed count.

\begin{figure}[h]
  \centering
  \includegraphics[width=0.98\linewidth]{fig3_pe_dominant_scaling.pdf}
  \caption{PE-dominant evidence (visualization of Table~\ref{tab:pe-dominant}). Left: EVQ's closed-form allocation (zero learned parameters) closes more of the $128{\to}8\mathrm{K}$ extrapolation gap than the learned positional-operator baseline. Right: at $L{=}256$, the predicted $\tau^*{=}4.0$ is favored among tested settings, and EVQ4+YaRN outperforms Geo+YaRN across extrapolation ratios.}
  \label{fig:pe-dominant}
\end{figure}

\begin{table}[h]
\caption{$L{=}256$ YaRN composition study (125M, PPL). The predicted $\tau^*{=}4.0$ is the best of the tested settings, and EVQ4+YaRN outperforms Geo+YaRN at every extrapolation ratio. This supplements Table~\ref{tab:pe-dominant} in the main text.}
\label{tab:pe-yarn-l256}
\centering\small
\setlength{\tabcolsep}{6pt}
\begin{tabular}{@{}lcccc@{}}
\toprule
Method (\(L{=}256\)) & 4$\times$ & 8$\times$ & 16$\times$ & 32$\times$ \\
\midrule
Geo+YaRN   & 60.7 & 97.5 & 190.4 & 260.2 \\
EVQ2+YaRN  & 51.6 & 66.7 & 107.9 & 141.7 \\
EVQ4+YaRN  & \textbf{50.9} & \textbf{60.9} & \textbf{84.1} & \textbf{99.6} \\
\bottomrule
\end{tabular}
\end{table}

\paragraph{MLA: allocation matters more when channels are scarce.}
EVQ extends to MLA \citep{deepseekv2, deepseekv3}, where only 16 frequency channels carry positional information. On a 432M model (3-seed, 500M tokens, $L_{\mathrm{train}}{=}8192$), EVQ alone reduces $2\times$ extrapolation PPL by $31.1\%$ and outperforms Geo+YaRN at $16\mathrm{K}$. Preliminary evidence under thorough training (1B tokens, single seed) suggests that EVQ's raw advantage reverses ($+11.1\%$) while EVQ+YaRN after fine-tuning still wins ($-2.5\%$)---a $13.6\mathrm{pp}$ swing consistent with the frequency substrate remaining the binding constraint even when raw extrapolation saturates. Multi-seed replication is in progress. Full 3-seed standalone results are in Appendix~\ref{sec:mla-results}.

\paragraph{Post-training injection via LoRA (7B scale).}
As supporting post-hoc evidence, we inject EVQ-Cosh frequencies ($\tau{=}1.414$, the $\tau^*$ predicted for $d_{\mathrm{head}}{=}128$, $L{=}8192$) into LLaMA-3-8B-Instruct \citep{grattafiori2024llama3} and fine-tune with LoRA ($r{=}64{=}K$, $\alpha{=}128$) on LongAlign-10k for 300 steps.

\begin{table}[ht]
\caption{EVQ-Cosh LoRA on LLaMA-3-8B-Instruct ($r{=}64$, 300 steps, single seed). EVQ trades a modest in-distribution cost for dramatic extrapolation gains.}
\label{tab:lora-8b}
\centering
\small
\begin{tabular}{@{}l r r r@{}}
\toprule
Context & Base PPL & EVQ-LoRA PPL & $\Delta$ \\
\midrule
8K (in-dist.) & \textbf{7.42} & 9.63 & $+30\%$ \\
16K ($2\times$) & 176.3 & \textbf{21.5} & $-88\%$ \\
32K ($4\times$) & 1942.5 & \textbf{104.3} & $-95\%$ \\
\bottomrule
\end{tabular}
\end{table}

Table~\ref{tab:lora-8b} shows that EVQ-LoRA reduces extrapolation PPL by $8\times$ at $16\mathrm{K}$ and $19\times$ at $32\mathrm{K}$ with a modest in-distribution cost ($+30\%$ PPL at $8\mathrm{K}$), serving as supporting post-hoc evidence at 7B scale. The rank is set at $r{=}d_{\mathrm{head}}/2{=}64$. The phenomenological variational model in \S\ref{sec:lora-regime} predicts a rank-dependent viability threshold near $r \approx d_{\mathrm{head}}/2$ under our calibration: below the threshold, coupling stiffness appears to prevent frequency redistribution; above it, extrapolation recovers. RULER scores do not improve, consistent with LoRA capacity sufficing for positional adaptation but not instruction-following at extrapolated lengths.

\paragraph{Progressive training (single seed, preliminary).}
Three-stage progressive training (454M, $512{\to}1024{\to}2048$) is reported as supporting evidence.
The EVQ+YaRN PPL gap at $16\mathrm{K}$ widens from $-34.6\%$ to $-81.2\%$, driven by Geo+YaRN degradation while EVQ+YaRN remains stable. At Stage~3, functional context reaches $48\mathrm{K}$ ($24\times$ training length) at PPL~$2.63$. After Stage~2, EVQ \emph{raw} (PPL $11.2$ at $16\mathrm{K}$) surpasses EVQ+YaRN at Stage~1 ($16.8$), suggesting that continued training can partially internalize the allocation benefit. This is a single-seed observation; multi-seed replication is in progress.

\paragraph{Cross-modal confirmation: video DiT (supporting).}
In 129.6M DiT head-to-head experiments (2~seeds, shared training run), EVQ wins every metric in both seeds, eliminating optimizer and CUDA confounds. Once the dead-channel bottleneck is removed (base-1000), both $\tau{=}1.2$ and $\tau{=}1.5$ beat Geo by ${\approx}48\%$ on far-frame MSE. This is a supporting mechanism check rather than a standalone video contribution: the collision bottleneck identified by the theory appears across modalities. Full results are in \S\ref{sec:tau-correction}.

\paragraph{Signal gradient across abstraction layers.}
Across the diagnostics reported here, EVQ's effect generally attenuates through successive abstraction layers: raw PPL $-46\%$ (Table~\ref{tab:multiscale}) $\to$ Gold-Answer NLL $-30\%$ (Table~\ref{tab:quality-nll}) $\to$ passkey $+59\mathrm{pp}$ vs.\ Geo baseline (EVQ+YaRN $100\%$ vs.\ Geo $41\%$, 3~seeds) $\to$ QA accuracy $+2.2\mathrm{pp}$ (Table~\ref{tab:quality-nll}, single seed, near-random baseline). This pattern is consistent with the intervention being largest at the PE layer and attenuated through higher abstraction where model capacity becomes the binding constraint.

\paragraph{Capability preservation.}
Table~\ref{tab:capability} shows that EVQ is non-destructive when retrieval is absent from training and improves extrapolated retrieval by $+10$--$13\mathrm{pp}$ when retrieval is learned (454M, 3-seed).

\begin{table}[t]
\caption{Capability preservation and retrieval robustness (454M, 3-seed, per-document PPL scoring). The same model architecture as Table~\ref{tab:evq-yarn}; PPL differences arise from the scoring protocol (per-document mean vs.\ full-sequence). EVQ is non-destructive when retrieval is absent, improves extrapolated retrieval, and is more robust to denser passkey supervision. PPL@$16\mathrm{K}$ differs from Table~\ref{tab:evq-yarn} ($253.2$ vs.\ $262.0$) due to the scoring protocol: Table~\ref{tab:evq-yarn} uses full-sequence scoring while this table uses per-document mean; the relative ranking is unchanged.}
\label{tab:capability}
\centering
\small
\begin{tabular}{@{}l c c c@{}}
\toprule
Setting & Geo & EVQ & Interpretation \\
\midrule
Pure FineWeb-Edu passkey & 55.7\% & 56.7\% & non-destructive \\
10\% mix retrieval @2K & 100\% & 100\% & both learn task \\
10\% mix retrieval @4K & 58.7\% & 68.7\% & +10.0pp \\
10\% mix retrieval @8K & 40.7\% & 53.3\% & +12.7pp \\
10\% mix global retrieval & 66.7\% & 74.0\% & +7.3pp \\
10\% mix PPL@2K & 67.2 & 67.9 & +1.0\% \\
10\% mix PPL@8K & 161.9 & 150.3 & -7.2\% \\
10\% mix PPL@16K & 262.0 & 237.2 & -9.5\% \\
5\%\,$\rightarrow$\,10\% delta @8K & -13.3pp & -3.3pp & EVQ more robust \\
\bottomrule
\end{tabular}
\end{table}

